\begin{longtable}{TD}
\caption{Resource demand, occupancy, cost, and overhead terms.}
\label{tab:terms-resource-demand-occupancy-cost}\\

\toprule
Term & Definition \\
\midrule
\endfirsthead

\toprule
Term & Definition \\
\midrule
\endhead

\bottomrule
\endfoot

\termrow{resource demand}{
The expected or required active resource need of a task attempt along a
resource dimension. Resource demand is what the model estimates before or
during execution.
}

\termrow{resource usage}{
The observed actual use of a resource during execution. Resource usage is
measured after or during execution and is distinct from estimated demand.
}

\termrow{active resource usage}{
Actual use of active resources such as CPU, memory, disk, network, runtime
allocation, or accelerator resources. Active resource usage is distinct from
waiting and bounded-capacity occupancy.
}

\termrow{CPU time}{
The amount of CPU execution time consumed by a task attempt. CPU time is not
the same as wall-clock duration.
}

\termrow{CPU pressure}{
Pressure on CPU capacity in the worker, container, node, or co-located
environment. CPU pressure is an execution-context factor, not the same as task
CPU demand.
}

\termrow{memory peak}{
The maximum memory used by a task attempt during execution. Memory peak is a
resource-usage or resource-demand dimension.
}

\termrow{memory pressure}{
Pressure on available memory in the worker, container, node, or runtime
environment. Memory pressure is an execution-context factor.
}

\termrow{allocation behavior}{
The allocation volume, allocation rate, allocation pattern, or
allocation-induced runtime effect of a task attempt. Allocation behavior may
influence memory pressure, garbage collection, and CPU time.
}

\termrow{disk I/O}{
Disk reads, writes, operations, bytes, or I/O wait associated with a task
attempt. Disk I/O demand is distinct from disk pressure.
}

\termrow{network I/O}{
Network bytes, packets, requests, or transfer volume associated with a task
attempt. Network I/O is distinct from network instability.
}

\termrow{dependency wait time}{
Time spent waiting for an external dependency, such as a database, API, remote
service, or storage system. Dependency wait time can increase worker occupancy
without increasing CPU time.
}

\termrow{worker occupancy}{
Bounded execution capacity held by a task attempt over time. Worker occupancy
is a first-class cost component in queue-worker systems.
}

\termrow{worker-slot time}{
The time during which a task attempt holds a worker execution slot. This is the
core worker-occupancy dimension.
}

\termrow{lease time}{
The time during which a task lease or reservation is held. Lease time may
differ from worker-slot time.
}

\termrow{connection hold time}{
The time during which a task attempt holds a connection to a dependency,
database, broker, or external service.
}

\termrow{lock hold time}{
The time during which a task attempt holds a lock. Lock hold time can affect
other tasks and may contribute to contention.
}

\termrow{lock wait time}{
The time spent waiting to acquire a lock. Lock wait time is distinct from lock
hold time.
}

\termrow{semaphore hold time}{
The time during which a task attempt holds a semaphore, concurrency token, or
bounded execution permit.
}

\termrow{tenant concurrency time}{
The time during which a task attempt consumes a tenant-level concurrency slot.
This dimension is relevant to fairness and isolation.
}

\termrow{wall-clock duration}{
Elapsed time from attempt start to attempt finish. Wall-clock duration is often
related to worker-slot time, but the two are not necessarily identical.
}

\termrow{latency}{
Delay observed by a caller, queue, worker, dependency, or system boundary.
Latency is an outcome or symptom, not resource demand by itself.
}

\termrow{service time}{
The time a task attempt spends being served or executed, excluding queue wait
time. The exact boundary must be defined for the queue-worker system being
analyzed.
}

\termrow{queue wait time}{
The time between enqueue and dequeue or execution start. Queue wait time is
distinct from worker-slot time.
}

\termrow{task cost}{
An umbrella term for resource demand, worker occupancy, and execution overhead
associated with a task attempt. Operational risk is not part of task cost.
}

\termrow{execution overhead}{
Overhead inherent to executing a task attempt, such as serialization, input
fetching, setup, acknowledgement, commit, cleanup, or retry bookkeeping.
}

\termrow{profiling overhead}{
Overhead caused by profiling and instrumentation, including tracing, metrics,
runtime profiling, high-cardinality labels, sampling decisions, and telemetry
export.
}

\termrow{telemetry overhead}{
CPU, memory, allocation, network, storage, or cardinality cost introduced by
telemetry collection and export.
}

\end{longtable}
